% ============================================================================
\section*{记号与约定}
\label{sec:notation}
\addcontentsline{toc}{section}{记号与约定}
% ============================================================================

本节列出全书常用的数学记号。初次阅读时可以跳过，遇到不熟悉的符号再回来查阅。

\subsection*{数集与基本结构}

\begin{longtable}{@{}p{3.5cm}p{10cm}@{}}
\toprule
\textbf{记号} & \textbf{含义} \\
\midrule
$\Z, \Q, \R, \C$ &
  整数、有理数、实数、复数。例如 $\Z = \{\ldots, -2, -1, 0, 1, 2, \ldots\}$。 \\[4pt]
$\Z/N\Z$ &
  模 $N$ 的整数（剩余类环）。元素为 $\{0, 1, 2, \ldots, N-1\}$，加法和乘法都"模 $N$"。例如 $\Z/5\Z$ 中 $3 + 4 = 2$。 \\[4pt]
$(\Z/N\Z)^\times$ &
  $\Z/N\Z$ 的\textbf{可逆元群}（unit group），即与 $N$ 互素的剩余类。例如 $(\Z/8\Z)^\times = \{1, 3, 5, 7\}$，共 $\varphi(8) = 4$ 个元素。 \\[4pt]
$\varphi(N)$ &
  \textbf{Euler $\varphi$-函数}：$1$ 到 $N$ 中与 $N$ 互素的整数个数。即 $\#(\Z/N\Z)^\times$。 \\[4pt]
$\F_p$ &
  含 $p$ 个元素的有限域（$p$ 为素数），即 $\Z/p\Z$。 \\[4pt]
$\HH$ &
  \textbf{上半平面}：$\HH = \{\tau \in \C \mid \im(\tau) > 0\}$。可以想象为复平面中实轴上方的半个平面。 \\[4pt]
$\hat\C = \C \cup \{\infty\}$ &
  \textbf{黎曼球面}：在复平面上加一个"无穷远点"。 \\
\bottomrule
\end{longtable}

\subsection*{矩阵群}

\begin{longtable}{@{}p{3.5cm}p{10cm}@{}}
\toprule
\textbf{记号} & \textbf{含义} \\
\midrule
$\GL_2(R)$ &
  \textbf{一般线性群}（General Linear group）：$R$ 上所有\textbf{可逆}的 $2\times 2$ 矩阵。"可逆"意味着行列式 $\det \neq 0$。当 $R = \Z$ 时，要求 $\det = \pm 1$。 \\[4pt]
$\SL_2(R)$ &
  \textbf{特殊线性群}（Special Linear group）：$R$ 上所有\textbf{行列式为 $1$} 的 $2\times 2$ 矩阵。例如 $\begin{pmatrix} 3 & 1 \\ 5 & 2 \end{pmatrix} \in \SL_2(\Z)$ 因为 $3 \cdot 2 - 1 \cdot 5 = 1$。字母 S 代表"Special"，意思是加了行列式$=1$的特殊约束。 \\[4pt]
$\PSL_2(R)$ &
  \textbf{射影特殊线性群}：$\PSL_2(R) = \SL_2(R) / \{\pm I\}$。就是在 $\SL_2$ 中把矩阵 $A$ 和 $-A$ 视为同一个元素。这是因为 $A$ 和 $-A$ 对上半平面的作用完全相同。 \\[4pt]
$I$ &
  单位矩阵 $\begin{pmatrix} 1 & 0 \\ 0 & 1 \end{pmatrix}$。 \\[4pt]
$-I$ &
  $\begin{pmatrix} -1 & 0 \\ 0 & -1 \end{pmatrix}$。在 $\SL_2(\Z)$ 中 $-I \neq I$，但它们对 $\HH$ 的作用相同。 \\
\bottomrule
\end{longtable}

\subsection*{同余子群}

\begin{longtable}{@{}p{3.5cm}p{10cm}@{}}
\toprule
\textbf{记号} & \textbf{含义} \\
\midrule
$\Gamma(N)$ &
  \textbf{主同余子群}：$\SL_2(\Z)$ 中所有模 $N$ 等于单位矩阵的矩阵。最"小"（最多约束）的同余子群。 \\[4pt]
$\GammaI(N)$ &
  $\SL_2(\Z)$ 中满足 $a \equiv d \equiv 1$，$c \equiv 0 \pmod{N}$ 的矩阵（$b$ 任意）。约束比 $\Gamma(N)$ 少一条。 \\[4pt]
$\GammaO(N)$ &
  $\SL_2(\Z)$ 中满足 $c \equiv 0 \pmod{N}$ 的矩阵。只有一条约束，所以是最"大"的。包含关系：$\Gamma(N) \subset \GammaI(N) \subset \GammaO(N) \subset \SL_2(\Z)$。 \\
\bottomrule
\end{longtable}

\subsection*{群论与代数记号}

\begin{longtable}{@{}p{3.5cm}p{10cm}@{}}
\toprule
\textbf{记号} & \textbf{含义} \\
\midrule
$H \leq G$ &
  $H$ 是 $G$ 的\textbf{子群}（subgroup），即 $H \subset G$ 且 $H$ 自身构成一个群。 \\[4pt]
$H \trianglelefteq G$ &
  $H$ 是 $G$ 的\textbf{正规子群}（normal subgroup），即对所有 $g \in G$ 都有 $gHg^{-1} = H$。正规子群可以用来做商群。 \\[4pt]
$[G : H]$ &
  $H$ 在 $G$ 中的\textbf{指标}（index），即陪集的个数。直觉："$G$ 比 $H$ 大多少倍"。 \\[4pt]
$G \backslash X$ &
  群 $G$ 作用在 $X$ 上的\textbf{商空间}（quotient）：把同一个轨道中的点全部"粘合"成一个点。$\SL_2(\Z) \backslash \HH$ 就是把上半平面中被 $\SL_2(\Z)$ 连接的点合并。 \\[4pt]
$G_x$ &
  $x$ 的\textbf{稳定子}（stabilizer）：所有把 $x$ 固定不动的群元素。即 $\{g \in G \mid g \cdot x = x\}$。 \\[4pt]
$G \cdot x$ &
  $x$ 的\textbf{轨道}（orbit）：$x$ 被群 $G$ 的所有元素映到的点的集合。 \\[4pt]
$\cong$ &
  \textbf{同构}（isomorphism）：两个数学结构"本质上一样"。例如 $\Z/2\Z \cong \{+1, -1\}$（乘法群）。 \\[4pt]
$\ker(\varphi)$ &
  映射 $\varphi$ 的\textbf{核}（kernel）：被映到单位元的所有元素。例如 $\SL_2(\Z) \to \SL_2(\Z/N\Z)$ 的核就是 $\Gamma(N)$。 \\
\bottomrule
\end{longtable}

\subsection*{模形式记号}

\begin{longtable}{@{}p{3.5cm}p{10cm}@{}}
\toprule
\textbf{记号} & \textbf{含义} \\
\midrule
$\Mk_k(\Gamma)$ &
  $\Gamma$ 上权 $k$ 的\textbf{模形式空间}。这是一个有限维向量空间。 \\[4pt]
$\Sk_k(\Gamma)$ &
  $\Gamma$ 上权 $k$ 的\textbf{尖点形式空间}。是 $\Mk_k(\Gamma)$ 的子空间。 \\[4pt]
$q$ &
  $q = e^{2\pi i \tau}$。当 $\tau \in \HH$ 时 $|q| < 1$。模形式的 Fourier 展开写成 $q$ 的幂级数。 \\[4pt]
$f(\tau) = \sum a_n q^n$ &
  模形式的 \textbf{$q$-展开}。系数 $a_n$ 编码了丰富的算术信息。 \\
\bottomrule
\end{longtable}
